% ATTEMPT_STATUS: unresolved
% ATTEMPT_ORIGIN: new research, 2026-10-01
% TOP_PROBLEM: {"id": 31, "title": "The Graceful Tree Conjecture", "category_id": 3, "category": {"id": 3, "name": "graph_theory", "display_name": "Graph Theory", "description": "Problems involving graphs, networks, and their properties.", "slug": "graph-theory", "order_index": 3, "created_at": "2026-07-31T15:26:25.671Z"}, "status": "open"}
% FIRST_POSTED: 2026-10-01
% Imported from: unsolvedmath_additional_500.tex; UnsolvedMath ID 31
% !TeX program = lualatex
% Compile twice: lualatex 31.tex
% Self-contained: no external bibliography, images, or \input files.
% Numeric section labels and visible numbers are original UnsolvedMath IDs.
\documentclass[11pt,a4paper]{article}
\usepackage[margin=24mm,headheight=14pt]{geometry}
\usepackage{fontspec}
\setmainfont{Latin Modern Roman}
\setsansfont{Latin Modern Sans}
\setmonofont{Latin Modern Mono}
\usepackage{amsmath,amssymb,mathtools}
\usepackage{unicode-math}
\setmathfont{Latin Modern Math}
\usepackage{microtype}
\usepackage{enumitem,xurl,needspace}
\usepackage[dvipsnames]{xcolor}
\usepackage{hyperref}
\hypersetup{unicode=true,colorlinks=true,linkcolor=MidnightBlue,urlcolor=MidnightBlue,
 pdftitle={UnsolvedMath Problem 31},pdfauthor={Source-annotated compilation},bookmarksnumbered=true}
\usepackage{bookmark,tocloft,titlesec,fancyhdr}
\setlength{\cftsecnumwidth}{4.2em}
\cftsetpnumwidth{2.5em}
\cftsetrmarg{4em}
\titleformat{\section}{\Large\bfseries\raggedright}{\thesection}{0.7em}{}
\pagestyle{fancy}\fancyhf{}
\fancyhead[L]{\small\sffamily UnsolvedMath research attempt}
\fancyhead[R]{\small Original catalogue IDs}
\fancyfoot[C]{\thepage}
\setlength{\parindent}{0pt}
\setlength{\parskip}{5pt plus 1pt}
\setlength{\emergencystretch}{3em}
\setcounter{tocdepth}{1}\setcounter{secnumdepth}{1}
\setlist[itemize]{leftmargin=3em,itemsep=4pt,topsep=4pt,parsep=0pt}
\allowdisplaybreaks
\newcommand{\N}{\mathbb{N}}
\newcommand{\Z}{\mathbb{Z}}
\newcommand{\Q}{\mathbb{Q}}
\newcommand{\R}{\mathbb{R}}
\newcommand{\C}{\mathbb{C}}
\newcommand{\F}{\mathbb{F}}
\newcommand{\A}{\mathbb{A}}
\newcommand{\PP}{\mathbb{P}}
\newcommand{\E}{\mathbb{E}}
\newcommand{\eps}{\varepsilon}
\newcommand{\dd}{\,\mathrm d}
\newcommand{\sref}[2]{\hyperlink{source:#1:#2}{[S#2]}}
\newif\ifshowcatalogue
\showcataloguetrue
\newcommand{\cataloguescope}[1]{\ifshowcatalogue
 \paragraph{Statement recorded in the supplied source.}
 #1\fi}
\begin{document}
% UnsolvedMath ID 31; source catalogue code GT-003

\setcounter{section}{30}

\section{The Graceful Tree Conjecture}\label{31}

{\small\sffamily UnsolvedMath ID 31 · Graph Theory}

\subsection{Definitions and mathematical statement}

\paragraph{Shared notation.}Unless a section says otherwise, $\N=\{1,2,\ldots\}$,
$\Z$ is the ring of integers, $\Q,\R,\C$ are the rational, real, and complex
fields, $[N]=\{1,\ldots,\lfloor N\rfloor\}$, and $\log$ is the natural
logarithm. For real functions of a parameter tending to infinity,
$f\ll g$ and $f=O(g)$ mean $|f|\leq Cg$ eventually for some fixed $C$;
$f\gg g$ reverses this comparison; $f=o(g)$ means $f/g\to0$;
$f\sim g$ means $f/g\to1$; and $f\asymp g$ means both order bounds.
A subscript on $O$ or $\ll$ specifies allowed constant dependence.
The expression $n^{o(1)}$ in an upper bound means growth smaller than
$n^\epsilon$ for each fixed $\epsilon>0$. Local definitions govern any
exception, finite-field convention, or use of the same symbol for another
object. An asymptotic claim is not automatically an effective algorithm.



A tree is a finite connected graph with no cycle. If it has $m$ edges, it has $m+1$ vertices. A graceful labelling is a bijection $f:V(T)\to\{0,\ldots,m\}$ for which the multiset $\{|f(u)-f(v)|:uv\in E(T)\}$ is exactly $\{1,\ldots,m\}$. Both the vertex labels and the induced edge differences are required to be distinct. \sref{31}{1} \sref{31}{2}

\paragraph{Statement recorded in the supplied source.}
Every tree can be gracefully labeled: vertices can be assigned distinct labels from $\{0, 1, \ldots, |E|\}$ such that edge labels (absolute differences) are all distinct. \sref{31}{1}

\subsection{Short English statement}

Can every finite tree be numbered so that its edge differences use each positive integer up to the number of edges exactly once?

\subsection{Research attempt}
\textbf{Status: UNRESOLVED.} We construct graceful labelings for paths and double stars, prove an extension criterion with a prescribed zero label, and exhibit the obstruction to turning that criterion into an induction for all trees. We also derive an exact finite search formulation and a parity check for any proposed labeling. No new general solution or new family theorem is claimed.

\paragraph{Source check and strategy.}
The inherited web formulation [S2] returned a 404 response when consulted on 1 October 2026. Gallian's primary survey [E1] gives an accessible reference for graceful graph labeling; only its framework is used here. We preserve the canonical statement. The proof attempt is constructive: arrange edge differences consecutively, then test whether a leaf can always be inserted into a graceful labeling of a smaller tree.

\paragraph{Paths: an explicit alternating construction.}
Let $v_0,\ldots,v_m$ be the vertices in path order. Define
\[
f(v_{2j})=j,\qquad f(v_{2j+1})=m-j
\]
whenever the subscripts lie between zero and $m$. The even-position labels increase from zero and the odd-position labels decrease from $m$. They form disjoint intervals whose union is $\{0,\ldots,m\}$, so $f$ is bijective. The edge from position $2j$ to $2j+1$ has difference $m-2j$; the next edge has difference $m-2j-1$. In path order the edge differences are therefore $m,m-1,\ldots,1$.

This includes $m=0$, where the single vertex receives zero and the edge condition is empty. It also places zero at one end of every path. The formula supplies a chosen labeling, not a license to put zero at an arbitrary prescribed vertex by an arbitrary permutation of labels.

\paragraph{Double stars: a second complete family.}
A double star has adjacent centres $u,v$, with $a$ leaves at $u$ and $b$ leaves at $v$, hence $m=a+b+1$ edges. Give $u$ label zero and $v$ label $b+1$. Give the $b$ leaves at $v$ the labels $1,\ldots,b$, and the $a$ leaves at $u$ the labels $b+2,\ldots,m$. These labels partition $\{0,\ldots,m\}$. The $v$-leaf differences are $1,\ldots,b$, the central difference is $b+1$, and the $u$-leaf differences are $b+2,\ldots,m$. Every desired difference occurs once.

When $a$ or $b$ is zero, the corresponding interval is empty and the same proof applies. The construction consequently contains stars and the single edge without additional exceptions. It relies on the two-level incidence pattern; when an attached branch has its own branches, assigning a consecutive interval at one centre no longer controls all internal differences.

\paragraph{The exact leaf-extension lemma.}
Suppose a tree $T$ has $m$ edges and a graceful labeling $f$ with $f(v)=0$. Attach a new leaf $w$ to $v$ and put $f(w)=m+1$, leaving old labels unchanged. The old edge differences remain $1,\ldots,m$, and the new edge has difference $m+1$. This proves gracefulness of the enlarged tree.

Conversely, if an extension leaves every old label unchanged and uses the only unused label $m+1$ on $w$, it is graceful only if $f(v)=0$: the old differences already exhaust $1,\ldots,m$, so the new difference must be $m+1$, forcing $|(m+1)-f(v)|=m+1$. Since $0\le f(v)\le m$, this is equivalent to $f(v)=0$.

Complementation $f\mapsto m-f$ preserves every difference and exchanges labels zero and $m$. It does not allow an arbitrary vertex to become zero. Similarly, permuting labels arbitrarily does not preserve edge differences. Thus induction by deleting a leaf requires a stronger rooted assertion at its attachment vertex, not merely the ordinary induction hypothesis that the smaller tree is graceful.

\paragraph{A failed extension of a specific graceful labeling.}
Take the three-vertex path $a-b-c$ with labels $(0,2,1)$. Its differences are two and one, so it is graceful. Attach a new leaf $w$ to $c$. Keeping the old labels and assigning three to $w$ creates difference two again; the required difference three is absent. Complementing the old labeling gives $(2,0,1)$ and still leaves label one at $c$, so this second unchanged-label extension also fails.

The enlarged tree is a four-vertex path and is graceful, for instance with labels $(0,3,1,2)$ in the order $a,b,c,w$. Thus the failure is not a counterexample to the conjecture; it proves that a successful induction may need to change multiple old labels. The usual induction hypothesis provides some graceful labeling, without the additional root condition needed by the extension lemma. We have not established a general theorem supplying that condition or a universally valid relabeling operation.

\paragraph{A useful certificate identity.}
For any graceful labeling,
\[
\sum_{uv\in E(T)}|f(u)-f(v)|=\frac{m(m+1)}2.
\]
Reducing modulo two and using $|r-s|\equiv r+s\pmod2$ yields
\[
\sum_{v\in V(T)}\deg(v)f(v)\equiv \frac{m(m+1)}2\pmod2.
\]
Only odd-degree vertices contribute on the left. This is a necessary checksum for a candidate. It is not sufficient: repeated differences can have the correct sum or parity. A sound verifier must test bijectivity of vertex labels and equality of the entire edge-difference multiset to $\{1,\ldots,m\}$.

\paragraph{Finite search and final gap.}
One exact formulation has integer variables $x_v\in\{0,\ldots,m\}$, pairwise distinct, and edge variables $y_{uv}=|x_u-x_v|\in\{1,\ldots,m\}$, also pairwise distinct. Since the numbers of variables match their allowed label sets, these constraints are equivalent to gracefulness. Enumerating vertex permutations is complete but factorial in $m+1$. Complementation permits a symmetry restriction on one fixed edge, such as $x_u<x_v$, without losing all representatives; prescribing a root label zero lacks that justification.

The constructions and identities above were checked algebraically, with no claimed exhaustive enumeration. The exact unresolved step is a labeling construction for every tree, or a replacement extension lemma that can relabel an existing tree while maintaining all old differences. The family formulas provide positive test cases for such a method. A finite search finding labelings for small trees would not prove that an arbitrary attachment vertex can always meet the induction requirement.

\subsection{Sources}

{\small

\begin{itemize}
\item[E1] J. A. Gallian, \emph{A Dynamic Survey of Graph Labeling}, 27th edition, 15 November 2024. Source consulted 1 October 2026. \url{https://www.combinatorics.org/files/Surveys/ds6/ds6v27-2024.pdf}.


\item[\hypertarget{source:31:1}{S1}] ULAM.AI, \emph{UnsolvedMath}, supplied \texttt{problems.json}, record 31 (GT-003), statement and background. The embedded literature-review date is 2026-08-17; that is the archive’s review date, not a fresh certification. Dataset landing page: \url{https://huggingface.co/datasets/ulamai/UnsolvedMath}.

\item[\hypertarget{source:31:2}{S2}] Graceful Tree Conjecture, Graph-theory open problems tracker. \url{https://mlelarge.github.io/graph-conjectures/op/graceful_tree_conjecture/}. Bibliographic information imported from the supplied record.

\end{itemize}

}

\label{end:31}
\end{document}
